% Generated draft: requires mathematical and source review.
\documentclass[11pt]{article}
\usepackage[margin=25mm]{geometry}
\usepackage{fontspec}
\setmainfont{Latin Modern Roman}
\usepackage{amsmath,amssymb,mathtools,unicode-math}
\setmathfont{Latin Modern Math}
\usepackage{enumitem,needspace,xurl,hyperref}
\setlength{\parindent}{0pt}
\setlength{\parskip}{5pt}
\setlength{\emergencystretch}{3em}
\newcommand{\N}{\mathbb{N}}
\newcommand{\Z}{\mathbb{Z}}
\newcommand{\Q}{\mathbb{Q}}
\newcommand{\R}{\mathbb{R}}
\newcommand{\C}{\mathbb{C}}
\newcommand{\F}{\mathbb{F}}
\newcommand{\A}{\mathbb{A}}
\newcommand{\PP}{\mathbb{P}}
\newcommand{\E}{\mathbb{E}}
\newcommand{\eps}{\varepsilon}
\newcommand{\dd}{\,\mathrm{d}}
\begin{document}
\section*{UnsolvedMath ID: 30000832}
\textbf{Source problem code:} catalogue entry
\subsection{Definitions and mathematical statement}
Let \(S\) be a real algebraic surface with isolated singularities, and let \(p\) be one of those singularities. Write \(S_{\mathrm{reg}}\) for its smooth locus. The question asks whether the normal-form metrics associated with such a general surface yield analogues of the higher-order asymptotic and regularity results for the exponential map that were proved for the modeled quasihomogeneous singularities, after a suitable resolution or blow-up.

For a specified Riemannian metric \(g\) on \(S_{\mathrm{reg}}\), the ordinary exponential map at a regular point \(q\) is \(\exp_q(v)=\gamma_{q,v}(1)\), whenever the geodesic with initial data \((q,v)\) exists up to time one. The source concerns behavior near an isolated singularity, but does not specify whether the map being expanded is based at regular points approaching \(p\), is an extension of this map, or is defined from geodesics emerging from \(p\). An ordinary tangent-space exponential map at a singular point is not silently assumed.

A resolution or blow-up is intended to replace the singular geometry by a space on which the relevant metric or map can be analyzed near the preimage of \(p\). The record gives neither that construction nor the normal-form metric, model theorem, expansion variable, expansion order, error bounds, or regularity class. In particular, it does not assert that the lifted metric becomes a nondegenerate smooth Riemannian metric. Accordingly the source supplies a generalization question, rather than a complete quantified asymptotic theorem. A precise extension must first identify the model results and then specify which of their conclusions are to hold for which general isolated real algebraic surface singularities.
\subsection{Short English statement}
Can the higher-order exponential-map asymptotics and regularity obtained for modeled quasihomogeneous singularities be extended to general real algebraic surfaces with isolated singularities, using normal-form metrics and a suitable resolution or blow-up? The model conclusions and geometric constructions are not reproduced.
\subsection{Sources}
\begin{itemize}
\item Local UnsolvedMath registry, permanent ID \(30000832\), source code catalogue entry, ``Exponential-Map Asymptotics Near Algebraic Surface Singularities.'' No individual proposer is supplied.
\item \emph{Analysis and Geometric Singularities} (2008), Oberwolfach Reports, DOI 10.4171/owr/2007/41, pp. 2361--2442: \url{https://doi.org/10.4171/owr/2007/41}. This is the sole supplied external reference.
\item Supplied literature triage, checked 2026-08-21: status ``open.'' It calls the extraction an incomplete conjectural extension of the quasihomogeneous analysis and reports no located published theorem establishing the extension. This dated assessment derives from the supplied OWR curation; time-sensitive status and cited claims require independent verification.
\end{itemize}
\end{document}
